% Secao 9 - Implicacoes de mensuracao e politica. Descritivo.

\section{Implicações de mensuração e política}
\label{sec:implicacoes}

A escolha da régua tem consequências concretas para quem um programa de acesso
alcança. Sob o conceito de distância, 1.373 municípios caem no quartil superior
e seriam classificados como desertos; sob o conceito de fluxo, outros 1.373 ---
mas não os mesmos. Os dois conjuntos coincidem em apenas 778 municípios (o
quadrante de deserto pleno) e discordam sobre 1.190.% src: profiles - flow_only/distance_only 595 cada, both 778
O número total de desertos quase não muda com a régua; a \emph{identidade} deles
muda por completo.

A assimetria importa porque tem direção. Migrar de quilômetros para fluxo
revelado faria 595 municípios --- cerca de 9,2 milhões de pessoas --- passarem a
qualificar como desertos pela primeira vez: são os municípios perto-mas-isolados,
hoje invisíveis a uma política de mapa. No sentido inverso, 595 municípios
(cerca de 8,1 milhões de habitantes), concentrados na Amazônia de viagem
fluvial, deixariam de ser sinalizados, porque seu atendimento efetivo é local
apesar da distância no mapa. Uma política calibrada por distância, portanto, não
erra ``de menos'' ou ``de mais''; erra de \emph{lugar} --- destina atenção a
municípios já servidos localmente enquanto deixa de fora aqueles cujo hospital
próximo não resolve a necessidade que exige complexidade.

O custo de corrigir é baixo. O burden de fluxo se constrói inteiramente a partir
de registros administrativos já existentes --- as internações do SIH ---, sem
necessidade de novas coletas ou pesquisas de campo. Um regulador que hoje aloca
leitos, equipes ou transporte sanitário por distância pode recomputar a mesma
geografia por fluxo revelado a custo essencialmente nulo, e checar quais de seus
desertos de mapa são reais e quais são artefato da régua.

Uma ressalva disciplina a leitura. O burden de fluxo é uma medida, não um
diagnóstico de carência: parte do fluxo para hubs distantes é referência
clínica apropriada --- casos de alta complexidade \emph{devem} concentrar-se em
centros de referência ---, não privação. Um burden alto sinaliza, portanto, um
município a inspecionar, não um cheque a emitir. O que o fluxo entrega é um mapa
melhor de onde olhar; a decisão de política continua exigindo o contexto local
e, para a pergunta sobre efeitos em saúde, o desenho causal da
Seção~\ref{sec:mortalidade}.
